\documentclass[aspectratio=169]{beamer}

% ------------------------------------------------
% THEME
% ------------------------------------------------
\usetheme{Madrid}
\usecolortheme{seahorse}
\usefonttheme{professionalfonts}
\setbeamertemplate{navigation symbols}{}

% ------------------------------------------------
% PACKAGES
% ------------------------------------------------
\usepackage{amsmath}
\usepackage{booktabs}
\usepackage{graphicx}
\usepackage{xcolor}

% ------------------------------------------------
% COLORS
% ------------------------------------------------
\definecolor{accent}{RGB}{0,92,184}

\setbeamercolor{structure}{fg=accent}
\setbeamercolor{title}{fg=white,bg=accent}

% ------------------------------------------------
% TITLE
% ------------------------------------------------
\title{Human Capital Theory: Applications to Education and Training}
\subtitle{Chapter 9}
\author{Prepared by Dr. Mir Ahasan Kabir}
\institute{York University}
\date{}

\begin{document}

%------------------------------------------------
\begin{frame}
\titlepage
\end{frame}

%------------------------------------------------
\begin{frame}{Learning Objectives}
\begin{itemize}
\item Understand how investments in education and training can be analyzed as an \textbf{economic investment decision}.
\item Recognize that the \textbf{opportunity cost of time} is often the largest cost of education.
\item Understand why comparing wages of educated and less educated workers can be misleading due to:
\begin{itemize}
\item ability bias
\item labour market signalling
\end{itemize}
\item Explain why the \textbf{return to education} is large and increasing in Canada.
\item Discuss when workers or employers should finance training programs.
\end{itemize}

\textbf{Explanation:}

Human capital theory treats education similarly to investments in physical capital. Individuals compare the \textbf{costs of schooling} with the \textbf{future benefits in higher earnings}.
\end{frame}

%------------------------------------------------
\begin{frame}{Main Questions}
\begin{itemize}
\item Why are more educated workers paid more than less educated workers?
\item What determines the market rate of return to education?
\item If education is so profitable, why doesn’t everyone pursue a PhD?
\item Do wages reflect the productivity gained through education or simply signal ability?
\item What roles do signalling and screening play in labour markets?
\end{itemize}

\textbf{Explanation}

These questions motivate the theory of human capital. Economists attempt to determine whether education actually \textbf{raises productivity} or simply acts as a \textbf{signal of ability}.
\end{frame}

%------------------------------------------------
\begin{frame}{Human Capital Theory}
\begin{block}{Definition}
\textbf{Human capital} refers to the knowledge, skills, and abilities acquired through education, training, and experience that increase productivity and earnings.
\end{block}

Key ideas:

\begin{itemize}
\item Investments in education involve \textbf{costs today}.
\item These investments generate \textbf{higher earnings in the future}.
\item Rational individuals invest if the \textbf{present value of benefits exceeds costs}.
\end{itemize}

\textbf{Example}

A student who attends university for four years sacrifices income during those years but expects higher wages over the rest of their career.
\end{frame}

%------------------------------------------------
\begin{frame}{Costs and Benefits of Human Capital Investment}

Important components include:

\begin{itemize}
\item \textbf{Opportunity cost}: forgone earnings while studying.
\item \textbf{Private costs and benefits}: costs and gains experienced by the individual.
\item \textbf{Social costs and benefits}: effects on society as a whole.
\item \textbf{Real costs}: resources used for education (teachers, buildings).
\item \textbf{Transfer costs}: financial transfers such as taxes or subsidies.
\end{itemize}

\textbf{Example}

Tuition is a private cost to students but may be subsidized by the government, shifting some costs to taxpayers.
\end{frame}

%------------------------------------------------
\begin{frame}{Private Investment in Education}

\begin{center}
%\includegraphics[width=.8\linewidth]{figure91}
\end{center}

\textbf{Explanation}

The diagram compares two lifetime income streams:

\begin{itemize}
\item income if the individual starts working immediately
\item income if the individual attends school first
\end{itemize}

Education typically reduces income early in life but increases income later.
\end{frame}

%------------------------------------------------
\begin{frame}{Age Earnings Profiles}

Empirical evidence shows that:

\begin{itemize}
\item Earnings increase with age and experience.
\item The increase occurs at a decreasing rate.
\item Workers with more education earn consistently higher wages.
\end{itemize}

\textbf{Example}

University graduates typically start earning later but surpass high school graduates within several years.
\end{frame}

%------------------------------------------------
\begin{frame}{Optimal Human Capital Investment}

Assumptions used in the model:

\begin{itemize}
\item Education has no direct consumption value.
\item Work hours are fixed.
\item Income streams are known with certainty.
\item Perfect capital markets allow borrowing and lending.
\end{itemize}

\textbf{Economic Idea}

Individuals choose the level of schooling that maximizes the \textbf{present value of lifetime earnings}.
\end{frame}

%------------------------------------------------
\begin{frame}{Net Present Value of Education}

For an 18-year-old deciding whether to attend college:

\[
PV = \sum_{t=18}^{T} \frac{Y}{(1+r)^t}
\]

where:

\begin{itemize}
\item $Y$ = annual income
\item $r$ = interest rate
\item $T$ = retirement age
\end{itemize}

Education is worthwhile if the \textbf{present value of higher future earnings exceeds the cost of schooling}.
\end{frame}

%------------------------------------------------
\begin{frame}{Marginal Costs and Benefits}

Cost of one more year of schooling:

\[
MC = Y + D
\]

where

\begin{itemize}
\item $Y$ = forgone earnings
\item $D$ = direct costs (tuition, books)
\end{itemize}

Benefit:

\[
MB = \Delta Y
\]

which represents the increase in future income.
\end{frame}

%------------------------------------------------
\begin{frame}{Optimal Education Rule}

Education should increase until:

\[
MB = MC
\]

or equivalently

\[
i = r
\]

where

\begin{itemize}
\item $i$ = internal rate of return to education
\item $r$ = market interest rate
\end{itemize}

\textbf{Interpretation}

If the return to education exceeds the interest rate, additional schooling is profitable.
\end{frame}

%------------------------------------------------
\begin{frame}{Implications of the Theory}

\begin{itemize}
\item Human capital investment occurs mainly early in life.
\item Individuals expecting shorter careers invest less.
\item Progressive income taxation may reduce incentives to invest in education.
\end{itemize}

\textbf{Example}

Workers nearing retirement rarely invest in long-term education programs.
\end{frame}

%------------------------------------------------
\begin{frame}{Factors Influencing Education}

Key determinants:

\begin{itemize}
\item \textbf{Income taxes}: progressive taxes reduce returns to education.
\item \textbf{Student loans}: reduce borrowing constraints and increase schooling.
\item \textbf{Government subsidies}: encourage investment in human capital.
\end{itemize}
\end{frame}

%------------------------------------------------
\begin{frame}{Worked Example: Returns to Education}

Mary can earn \$30,000 immediately.

Interest rate = 5\%.

If education raises income to \$33,000:

\[
i = \frac{\Delta Y}{Y+D}
\]

If direct costs are zero:

\[
i = \frac{3000}{30000} = 10\%
\]

Education is profitable because the return exceeds the interest rate.
\end{frame}

%------------------------------------------------
\begin{frame}{Education as a Filter}

\begin{block}{Signalling Theory}
Education may act as a signal of worker productivity rather than increasing productivity itself.
\end{block}

Employers use education to screen workers because:

\begin{itemize}
\item ability is difficult to observe directly
\item educational achievement may reveal ability
\end{itemize}
\end{frame}

%------------------------------------------------
\begin{frame}{Empirical Evidence: Education and Earnings}

Key patterns:

\begin{itemize}
\item Earnings rise with experience.
\item Workers with more education earn significantly more.
\item The earnings gap widens with age.
\end{itemize}

\begin{center}
%\includegraphics[width=.8\linewidth]{figure94}
\end{center}
\end{frame}

%------------------------------------------------
\begin{frame}{Human Capital Earnings Function}

Economists estimate returns to education using:

\[
\ln Y = \alpha + rS + \beta_1 EXP + \beta_2 EXP^2 + \epsilon
\]

Where:

\begin{itemize}
\item $Y$ = earnings
\item $S$ = years of schooling
\item $EXP$ = experience
\end{itemize}

The coefficient $r$ estimates the \textbf{return to education}.
\end{frame}

%------------------------------------------------
\begin{frame}{Ability Bias}

Estimating returns to education is complicated because:

\begin{itemize}
\item ability
\item motivation
\item perseverance
\end{itemize}

may affect both education and wages.

Researchers address this using:

\begin{itemize}
\item twin studies
\item natural experiments
\item compulsory schooling laws
\end{itemize}
\end{frame}

%------------------------------------------------
\begin{frame}{Returns to Education and Inequality}

Evidence suggests:

\begin{itemize}
\item returns to education have increased since the mid-1990s
\item higher returns contribute to rising wage inequality
\item technological change increases demand for skilled workers
\end{itemize}
\end{frame}

%------------------------------------------------
\begin{frame}{Training}

Two types of training:

\textbf{General training}

\begin{itemize}
\item transferable across firms
\item workers usually pay the cost
\end{itemize}

\textbf{Specific training}

\begin{itemize}
\item valuable only within one firm
\item firms often pay for the training
\end{itemize}
\end{frame}

%------------------------------------------------
\begin{frame}{Government Training Programs}

Government intervention may be justified when:

\begin{itemize}
\item private markets underinvest in training
\item information is imperfect
\item disadvantaged workers face borrowing constraints
\end{itemize}

Training subsidies may increase employment and earnings.
\end{frame}

%------------------------------------------------
\begin{frame}{Chapter Summary}

\begin{itemize}
\item Education is an investment in human capital.
\item Individuals choose schooling by comparing costs and benefits.
\item Education may increase productivity or signal ability.
\item Empirical evidence shows strong returns to education.
\item Training programs and government policies influence human capital investment.
\end{itemize}

\end{frame}

%------------------------------------------------
\begin{frame}
\centering
\Huge Thank You
\end{frame}

\end{document}